% !TeX root = ../教材题库_全部条目_转写工作稿.tex
% Source: XJ-B1-C05-EX-P190-Y0361
% 第5章 / 5.4 函数 y=A sin(ωx+φ) 的图象与性质 / 教材 p.190 / PDF 第197页
% OCR未校对：公式与文字均可能有误；未完成数学TeX转写；禁止直接发给学生
% 原题及相邻原解截图（核对用）：../题目原图/XJ-B1-C05-EX-P190-Y0361-01.webp
% 本文件只在转写完成后作为题干使用；OCR 条目尚不是可靠题干。
\begin{bankitem}{5.4 例题 例题2}{第5章 / 5.4 函数 y=A sin(\ensuremath{\omega}x+\ensuremath{\varphi}) 的图象与性质 / 教材 p.190 / PDF 第197页}{XJ-B1-C05-EX-P190-Y0361}{OCR未校对：公式与文字均可能有误；未完成数学TeX转写；禁止直接发给学生}
\begin{ocrdraft}
[教材p.190 / PDF第197页]
\par
0二0]在同一平面直角坐标系中画出3-艹“y=sin20一在[一2丌，
\par
2丌]上的图象，观察它们之间的关系，并说出这三个函数的周期、最大值、最小值、
\par
值域之间的关系．
\par
解通过“五点法”作出函数3，一“3，一sin2“3，一在区间[一2丌，
\par
2丌]上的简图，如图5．4一5．
\par
Y
\par
y=sin2x
\par
y=sln—X
\par
一2丌
\par
2訌
\par
2
\par
2
\par
2
\par
图5．4一5
\par
观察图5．4一5，可以发现：3，一sin2的图象可以由一sinx的图象上每一点
\par
0，sinx)的纵坐标不变、横坐标除以2（到3，轴的距离缩短到原来的砉）得到，因而
\par
y=sin2的值域、最大值、最小值都与3，一sin相同，周期缩短为
\par
2
\par
类似地，观察图5，4一5，可以发现：y一艹亠r的图象可以由y=sin的图象
\par
上每一点0，sin到的纵坐标不变、横坐标除以（到3，轴的距离放大到原来的2倍）
\par
得到．因而y=sin的值域、最大值、最小值都与y=sin相同，周期扩大为
\par
2丌\ensuremath{\cdot}2=4丌．
\par
一般地，对任意仞>0且仞1，函数3，一-的图象可以由y=sinx的图象
\par
上每一点的纵坐标不变、横坐标伸长（0<仞<1）或缩短\ensuremath{\text{\textperthousand}}>1）为原来的一得到．3-
\par
的值域为[一1，0，周期为
\end{ocrdraft}
\end{bankitem}
